\section*{2. ARCHITECTURAL CHARACTERISTICS}

\phantomsection
\addcontentsline{toc}{section}{2. ARCHITECTURAL CHARACTERISTICS}

\subsection*{2.1. Identifying Architecture Characteristics}

\phantomsection
\addcontentsline{toc}{subsection}{2.1. Identifying Architecture Characteristics}

We looked for characteristics in three places: the requirements in Chapter 1, the way TMC's factories actually operate, and facts that are not written in the assignment but hold for any gate system. A characteristic is kept only if it is non-functional, changes the structure of the system, and causes real damage to TMC when it is missed.

\Needspace{9\baselineskip}
\subsubsection*{2.1.1. Explicit Requirements}

\phantomsection
\addcontentsline{toc}{subsubsection}{2.1.1. Explicit Requirements}


Table 2.1 links each explicit requirement to the characteristic it creates and explains why it matters for the structure. In Table 2.1, characteristic names follow ISO/IEC 25010:2023, with sub-characteristics in brackets. Data integrity and privacy in this table, and labels such as Auditability later in the chapter, are our own terms, not names from the standard.

\begin{longtable}{|>{\raggedright\arraybackslash}p{87.5pt}|>{\raggedright\arraybackslash}p{110pt}|>{\raggedright\arraybackslash}p{212pt}|}
\hline
\multicolumn{1}{|>{\columncolor{tableheader}\centering\arraybackslash}p{87.5pt}|}{\textbf{Requirements}} & \multicolumn{1}{>{\columncolor{tableheader}\centering\arraybackslash}p{110pt}|}{\textbf{Characteristic}} & \multicolumn{1}{>{\columncolor{tableheader}\centering\arraybackslash}p{212pt}|}{\textbf{Architectural Significance}} \\
\hline
\endhead
\multicolumn{3}{c}{\rule{0pt}{20pt}\normalsize\textbf{Table 2.1. Architecture Characteristics Derived from Explicit Requirements}} \\
\endlastfoot
NFR-01; IC-02 & Performance efficiency (time behaviour) & The kiosk must show a result within 2 seconds at about 800 events per 30 minutes per factory. The turnstile call is included in this time, so the processing path at the gate must be short. \\
\hline
NFR-02; FR-03.3--FR-03.6 & Reliability (fault tolerance, recoverability); data integrity & No record may be lost during normal operation, and a resent scan must not be counted twice. The system must also distinguish between a gate that has opened and an employee who has actually passed. \\
\hline
NFR-03; FR-04.5; FR-04.7; FR-05.1; FR-05.5 & Security (accountability, authenticity) & Every operation must be limited by role and traceable to the user or device that performed it. \\
\hline
NFR-04; FR-01.2 & Security (confidentiality); privacy & Personal data appears in QR codes, kiosk screens, logs, and backups. Each of them should contain only the data that is needed. \\
\hline
IC-02; FR-02.2; I-01; I-03; I-05; I-07 & Compatibility (interoperability); reliability & The assignment does not define the controller signals for opening, passage, lane clear, and closing. A lost response must be treated as unknown, not as success or failure. \\*
\hline\noalign{\nobreak}
IC-01; FR-04.3; I-02 & Compatibility (interoperability); reliability & Data must reach HR/Payroll before noon on the first day of the month. A batch counts as delivered only when HR/Payroll acknowledges it. \\
\hline
\end{longtable}
\phantomsection\addcontentsline{lot}{table}{\protect\numberline{2.1}Architecture Characteristics Derived from Explicit Requirements}


\Needspace{9\baselineskip}
\subsubsection*{2.1.2. Domain Concerns}

\phantomsection
\addcontentsline{toc}{subsubsection}{2.1.2. Domain Concerns}


Several facts about how TMC's factories work also shape the characteristics:

\begin{itemize}[leftmargin=36pt, itemsep=2pt, topsep=4pt, parsep=2pt]

  \item \textbf{Shift changes:} About 800 people pass the gates of each factory within 30 minutes. A slow or unavailable gate quickly causes long queues, so performance and availability directly affect production.

  \item \textbf{Trust in attendance data:} HR staff and the HR/Payroll system use these records for payroll reconciliation. A lost or duplicated record may lead to incorrect pay, so data integrity must be maintained from the gate to the handover.

  \item \textbf{Accountability:} When a record is disputed, HR needs to know who changed a permission, who opened a gate in an emergency, or who approved an adjustment. Auditability is therefore required.

  \item \textbf{Role-based data access:} The six user groups have different responsibilities and data scopes. Security and privacy must limit both functions and data by role and scope.

  \item \textbf{Two sites in different provinces:} A failure of a shared component should not stop both factories, and a network failure at one factory should not affect the other. This affects availability and where components are deployed.

  \item \textbf{External signals:} The system depends on the controller and on HR/Payroll to report what happened. A missing or unclear signal needs a status check or reconciliation, which is an interoperability concern.

\end{itemize}

\Needspace{9\baselineskip}
\subsubsection*{2.1.3. Implicit Domain Knowledge}

\phantomsection
\addcontentsline{toc}{subsubsection}{2.1.3. Implicit Domain Knowledge}


Some facts are not written in the assignment but hold for any gate system of this kind:

\begin{itemize}[leftmargin=36pt, itemsep=2pt, topsep=4pt, parsep=2pt]

  \item \textbf{Bursty arrivals:} The rate of 0.44 events per second is only an average. Since most employees arrive shortly before the shift starts, we assume that half of the 800 events occur within 5 minutes: 400 / 300 s $\approx$ 1.33 events per second, about three times the average. These 400 events are completed entries and exits. Rejected and resubmitted scans come on top of them, so performance tests should add them to this rate instead of counting them inside it.

  \item \textbf{Independent failures:} Kiosks, network links, servers, databases, controllers, and HR/Payroll can each fail at a different time. The design must separate a failure before opening from an unknown opening result, a missing passage, and a failed close.

  \item \textbf{Resent requests:} When a response is lost, the kiosk may send the same scan again. The system must recognize the duplicate so that the gate is not opened twice and no second record is created.

  \item \textbf{Data outside the database:} Employee data also ends up in logs, error messages, and backups. Data protection has to cover these copies as well.

  \item \textbf{Changing rules:} Attendance rules, permission scopes, and API contracts will change during operation. Clear module boundaries keep such changes local, which supports maintainability.

\end{itemize}

\Needspace{9\baselineskip}
\subsection*{2.2. Characteristic Evaluation and Selection}

\phantomsection
\addcontentsline{toc}{subsection}{2.2. Characteristic Evaluation and Selection}


We rated each candidate High (H), Medium (M), or Low (L) on two scales. Importance asks what happens to TMC if the characteristic is missed. Difficulty asks how hard it is to achieve and check with two factories, two external systems, and a 2-second limit. A characteristic becomes a driver when its importance is High and it changes where components or deployments are split. A high difficulty tells us where more analysis is needed; it does not rank the drivers. Table 2.2 shows the result.

\begin{longtable}{|>{\raggedright\arraybackslash}p{36.5pt}|>{\raggedright\arraybackslash}p{158.5pt}|>{\raggedright\arraybackslash}p{63.5pt}|>{\raggedright\arraybackslash}p{63.5pt}|>{\raggedright\arraybackslash}p{62.5pt}|}
\hline
\multicolumn{1}{|>{\columncolor{tableheader}\centering\arraybackslash}p{36.5pt}|}{\textbf{ID}} & \multicolumn{1}{>{\columncolor{tableheader}\centering\arraybackslash}p{158.5pt}|}{\textbf{Characteristic}} & \multicolumn{1}{>{\columncolor{tableheader}\centering\arraybackslash}p{63.5pt}|}{\textbf{Importance}} & \multicolumn{1}{>{\columncolor{tableheader}\centering\arraybackslash}p{63.5pt}|}{\textbf{Difficulty}} & \multicolumn{1}{>{\columncolor{tableheader}\centering\arraybackslash}p{62.5pt}|}{\textbf{Driving}} \\
\hline
\endhead
\multicolumn{5}{c}{\rule{0pt}{20pt}\normalsize\textbf{Table 2.2. Evaluation and Selection of Architecture Characteristics}} \\
\endlastfoot
AC-01 & Reliability/Data Integrity & H & H & Yes \\
\hline
AC-02 & Availability & H & H & Yes \\
\hline
AC-03 & Performance & H & M & Yes \\
\hline
AC-04 & Security/Privacy & H & H & Yes \\
\hline
AC-05 & Auditability & H & M & Yes \\
\hline
AC-06 & Interoperability & H & M & Yes \\
\hline
- & Maintainability & M & M & No \\
\hline
- & Usability & M & L & No \\
\hline
- & Scalability/Elasticity & L & M & No \\*
\hline\noalign{\nobreak}
- & Portability & L & L & No \\
\hline
\end{longtable}
\phantomsection\addcontentsline{lot}{table}{\protect\numberline{2.2}Evaluation and Selection of Architecture Characteristics}


\Needspace{9\baselineskip}
\subsubsection*{2.2.1. Driving Characteristics}

\phantomsection
\addcontentsline{toc}{subsubsection}{2.2.1. Driving Characteristics}


Six characteristics are drivers, and each one shapes a different part of the structure:

\begin{enumerate}[leftmargin=36pt, itemsep=2pt, topsep=4pt, parsep=2pt]
  \item \textbf{AC-01 Reliability/Data Integrity} fixes the order of work at the gate: the accepted scan is stored before the opening command, and attendance is created only after passage. Every event that leaves the factory needs an identifier so that it is never counted twice.
  \item \textbf{AC-02 Availability} decides where gate work runs. A shift change must not stop because a factory's WAN link is down. A failure at the central site is different: new QR codes come only from there, so the automatic flow stops within about 60 seconds and guards switch to TMC's manual procedure.
  \item \textbf{AC-03 Performance} decides what stays on the synchronous path: the QR check, the access check, the durable write, and the controller call. Reports and the HR/Payroll handover are kept off this path.
  \item \textbf{AC-04 Security/Privacy} decides where login, permission checks, and data limits are enforced, and what may appear in QR codes, kiosk screens, and logs.
  \item \textbf{AC-05 Auditability} requires every gate, permission, review, and handover action to leave evidence where it happens.
  \item \textbf{AC-06 Interoperability} puts the controller and HR/Payroll behind their own adapters, because neither API is fully specified and a lost response has to be treated as unknown.
\end{enumerate}

\Needspace{6\baselineskip}
Table 2.3 gives a measurable success criterion for each driver. The values are targets and still have to be confirmed by testing.

{\small\renewcommand{\arraystretch}{1.08}
\begin{longtable}{|>{\raggedright\arraybackslash}p{70pt}|>{\raggedright\arraybackslash}p{135pt}|>{\raggedright\arraybackslash}p{215pt}|}
\hline
\multicolumn{1}{|>{\columncolor{tableheader}\centering\arraybackslash}p{70pt}|}{\textbf{Driver}} & \multicolumn{1}{>{\columncolor{tableheader}\centering\arraybackslash}p{135pt}|}{\textbf{Measure}} & \multicolumn{1}{>{\columncolor{tableheader}\centering\arraybackslash}p{215pt}|}{\textbf{Success criterion}} \\
\hline
\endhead
\multicolumn{3}{c}{\rule{0pt}{20pt}\normalsize\textbf{Table 2.3. Success Criteria for the Driving Characteristics}} \\
\endlastfoot
AC-01 Reliability / Data Integrity & Accepted scans, confirmed passages, duplicate scanIds, unresolved outcomes & Store ACCEPTED before openGate; create one automatic attendance transaction only for confirmed passage. Retrying the same scanId creates neither another command nor another transaction. Retain uncertain outcomes for reconciliation. \\
\hline
AC-02 Availability & Fault detection, recovery time, eligibility-data age & Report known gate status and preserve evidence when a component fails. After 60 seconds without confirmed updates, the factory may switch to degraded mode if TMC approves I-10, and it stops automatic acceptance 30 minutes after the last confirmed update, or earlier if another required condition fails. Detection and failover targets are proposed in Chapter 5. \\
\hline
AC-03 Performance & QR-read-to-kiosk-result time at 800 entry/exit events per 30 minutes per factory & Display a confirmed, durably stored opening result within 2 seconds (NFR-01). Valid scans that only receive UNDETERMINED count against this target; deliberately induced failures are measured separately. \\
\hline
AC-04 Security / Privacy & Unauthorized operations, invalid QR codes, exposed personal data & Reject requests outside the user's or device's scope, invalid or expired QR codes, and ineligible employees. Limit data in QR codes, displays, logs, and stores. Assisted attendance requires HR approval (I-04). \\
\hline
AC-05 Auditability & Coverage and queryability of action evidence & Record actor or device, action, target, time, outcome, and correlation identifiers for gate, permission, review, and handover actions. Authorized users can query evidence within their scope. \\*
\hline\noalign{\nobreak}
AC-06 Interoperability & Controller outcomes and HR/Payroll batch acknowledgements & Distinguish confirmed opening, passage, clearance, closure, and uncertain outcomes under I-01/I-05. Confirm batch receipt under I-02 and make required monthly data available before 12:00 noon on day one. \\
\hline
\end{longtable}
}
\phantomsection\addcontentsline{lot}{table}{\protect\numberline{2.3}Success Criteria for the Driving Characteristics}


The controller signals behind these criteria still have to be checked against the real API.

\Needspace{9\baselineskip}
\subsubsection*{2.2.2. Other Considered Characteristics}

\phantomsection
\addcontentsline{toc}{subsubsection}{2.2.2. Other Considered Characteristics}


We also considered the following characteristics but did not make them drivers:

\begin{itemize}[leftmargin=36pt, itemsep=2pt, topsep=4pt, parsep=2pt]

  \item \textbf{Maintainability (M/M):} The assignment asks us to show how SOLID supports maintenance and extension, but sets no specific change target that requires separate deployment units. Clear module interfaces support this goal; availability and the gate response time drive the deployment boundaries.

  \item \textbf{Usability (M/L):} Employees and guards must read results easily and tell a rejection from an unknown result. Simple screens and role-based interfaces are enough, so it does not shape the structure.

  \item \textbf{Scalability/Elasticity (L/M):} About 1,400 employees and roughly 1.3 events per second per factory at peak do not need automatic scaling. A second instance would be there to survive failures, not to carry more load. Running several instances would need shared scan state and duplicate control, so difficulty is Medium.

  \item \textbf{Portability (L/L):} No requirement asks for more than one deployment environment.

\end{itemize}

\Needspace{9\baselineskip}
\subsection*{2.3. Trade-off Analysis}

\phantomsection
\addcontentsline{toc}{subsection}{2.3. Trade-off Analysis}


Some characteristics pull in opposite directions. Table 2.4 lists the main conflicts and how we handle each one. We only weigh options that already meet the mandatory requirements.

{\renewcommand{\arraystretch}{1.2}
\begin{longtable}{|>{\raggedright\arraybackslash}p{93.5pt}|>{\raggedright\arraybackslash}p{138.5pt}|>{\raggedright\arraybackslash}p{177.5pt}|}
\hline
\multicolumn{1}{|>{\columncolor{tableheader}\centering\arraybackslash}p{93.5pt}|}{\textbf{Trade-off}} & \multicolumn{1}{>{\columncolor{tableheader}\centering\arraybackslash}p{138.5pt}|}{\textbf{Cause}} & \multicolumn{1}{>{\columncolor{tableheader}\centering\arraybackslash}p{177.5pt}|}{\textbf{Approach}} \\
\hline
\endhead
\multicolumn{3}{c}{\rule{0pt}{20pt}\normalsize\textbf{Table 2.4. Main Architectural Trade-offs}} \\
\endlastfoot
AC-04 Security vs AC-03 Performance & QR, status, and permission checks and the durable write with its local audit evidence all happen inside the 2 seconds. & Keep every mandatory check but avoid extra queries. Send audit records to Audit Log later. Use the factory's local access data only within the freshness limits of I-10. \\
\hline
AC-01 Reliability vs AC-03 Performance & Saving the accepted scan before opening, and the result before replying, adds time. & Keep both writes and the opening confirmation in the 2-second path. Measure passage-to-attendance and closing separately. \\
\hline
AC-06 Interoperability vs AC-01, AC-03 & Slow or unclear answers from the controller or HR/Payroll affect time and status. & Put each external system behind its own boundary. Keep the controller call inside the time budget and the HR/Payroll handover outside the scan flow (I-01--I-03). \\
\hline
AC-04 Privacy vs AC-05 Auditability & Tracing needs identifiers, times, and outcomes; privacy needs less data. & Store only what tracing needs. Raw QR tokens and secrets are never logged. Limit log access and set retention periods. \\
\hline
AC-04 Privacy vs reporting & Detailed reports help reconciliation but may show more than a user should see. & Limit reports by role and scope. Use aggregates where details are not needed. \\
\hline
AC-02 Availability vs AC-04 Security & If authorization data or links fail, opening anyway would rely on old permissions. & Authorization is never skipped to keep a gate running. Local data may be used during an outage only for a fixed time (I-10), with replay protection and a review flag on every decision made that way. \\*
\hline\noalign{\nobreak}
AC-01, AC-02 vs cost & Redundancy, backups, and recovery add development and operating cost. & Add only the mechanisms the requirements and risks justify, and state the cost of each. \\
\hline
\end{longtable}
}
\phantomsection\addcontentsline{lot}{table}{\protect\numberline{2.4}Main Architectural Trade-offs}


